% Step 2.1: geometry initialisation and tube-aligned Gaussians.
% Numbers from results/step2.1_geom_init_v60_fdk_run4.json (and run 1 for the static R2-Gaussian route).
% Bib keys: liu2025fourdrgs, zha2024r2gaussian, biguri2016tigre.

\paragraph{Geometry initialization.}
Because each view of a rotational DSA sweep is acquired at a different
time, a static reconstruction sees an inconsistent object. We therefore
initialize geometry from a ground-truth-free FDK volume of the training
views (Hann filter, TIGRE~\cite{biguri2016tigre}). The centerline is the
longest geodesic path of the skeleton of the thresholded volume, and the
lumen radius $R(s)$ at each arc length $s$ is the distance at which the
cross-sectional profile falls to half of its central value, taken as the
median over 16 directions. This half-maximum criterion is independent of
the intensity scale, which matters because each vessel segment is
opacified for a different fraction of the sweep. We place one Gaussian
every 0.5\,mm along the centerline, map it to its arc length $s_i$, align
its major axis with the local tangent, and set its cross-sectional scale
to $\sigma_i = R(s_i)/\sqrt{2}$ so that the lumen area is
$A(s_i) = 2\pi\sigma_i^2$. Scales use the bounded activation of
4DRGS~\cite{liu2025fourdrgs}, and a tube-alignment term
$\frac{1}{N}\sum_i \bigl(1 - (\mathbf{a}_i^\top \mathbf{t}_i)^2\bigr)$
keeps each major axis $\mathbf{a}_i$ along the tangent $\mathbf{t}_i$
during optimization. On five validation cases with 60 views, the
initial radius error is 0.05 to 0.12\,mm (mean 0.08\,mm), whereas a static
R$^2$-Gaussian~\cite{zha2024r2gaussian} fitted to the same dynamic views
drifts away from the vessel (0.20 to 0.42\,mm).
